\documentclass[11pt]{article}
\usepackage[margin=1in]{geometry}
\usepackage{amsmath,amssymb,amsthm,mathtools}
\usepackage{graphicx}
\usepackage{microtype}
\usepackage[hidelinks]{hyperref}

\newtheorem{theorem}{Theorem}
\newtheorem{lemma}{Lemma}
\newcommand{\R}{\mathbb{R}}
\newcommand{\E}{\mathcal{E}}
\newcommand{\A}{\mathcal{A}}

\title{An Exact Finite-$p$ Two-Radius Classification\\
for Horizontal Spherical Means on $H$-Type Groups}
\author{Candidate substantial partial result for arXiv:2109.14963}
\date{22 July 2026}

\begin{document}
\maketitle

\noindent\textbf{Status.}
Candidate partial result, likely valid, pending specialist review.  It settles
the source's two-radius question for every finite $p$ and leaves only the
nonexceptional $p=\infty$ endpoint.

\section{Source problem}

E. K. Narayanan, P. K. Sanjay, and K. T. Yasser ask in Remark~4.1 of
\cite{N_SY} whether a two-radius theorem holds above their one-radius
injectivity threshold $2m/(m-1)$, including $p=\infty$.

\begin{center}
\includegraphics[width=\linewidth]{figures/open_problem_crop.png}
\end{center}

The complete remark appears on physical page 19 of the arXiv PDF.  The part
addressed here is:
\begin{quote}
``It is a natural question whether a two radius theorem is true for functions
in $L^p(G)$ for $\frac{2m}{m-1}<p\leq\infty$.''
\end{quote}

\section{Definitions and spectral input}

Let $G$ be an $H$-type group whose Lie algebra splits orthogonally as
$\mathfrak v\oplus\mathfrak z$, with $\dim\mathfrak v=2n$ and
$\dim\mathfrak z=m\geq2$.  We write points as $(z,t)\in\R^{2n}\times\R^m$.
For $r>0$, let $\mu_r$ be normalized surface measure on
$\{(z,0):|z|=r\}$, and call $f\mapsto f*\mu_r$ the horizontal spherical
mean at radius $r$.

Let $L_k^{n-1}$ be the generalized Laguerre polynomial.  For $k\geq1$,
write its $k$ positive zeros as
\[
  0<\zeta_{k,1}<\cdots<\zeta_{k,k},
\]
and set
\[
 \E_n=
 \left\{\sqrt{\frac{\zeta_{k,i}}{\zeta_{k,j}}}:
 k\geq1,\ 1\leq i,j\leq k\right\}.
\]
This is a countable subset of $(0,\infty)$.

We use two results proved in the source paper \cite{N_SY}.  First, its
spectral projections $\A_k$ are bounded on $L^p(G)$ for
$1<p<\infty$.  Second, for $2\leq p<\infty$,
\begin{equation}\label{eq:abel}
 f=\lim_{\rho\uparrow1}\sum_{k=0}^{\infty}\rho^k\A_kf
 \quad\text{in }L^p(G).
\end{equation}
The proof of source Theorem~1.1 also establishes the following support fact
before it invokes its low-$p$ Fourier-support theorem.

\begin{lemma}[support consequence of the source proof]\label{lem:support}
Let $2\leq p<\infty$, $f\in L^p(G)$, and $u>0$.  If
$f*\mu_u=0$, then for every $k\geq0$ and almost every horizontal point $z$,
the distributional Fourier transform in $t$ of $\A_kf(z,\cdot)$ is supported
in
\[
 Z_{k,u}=\left\{a\in\R^m:
 L_k^{n-1}\!\left(\frac{|a|u^2}{2}\right)=0\right\}.
\]
\end{lemma}

\noindent The source states this inside the proof for
$p\leq2m/(m-1)$, but all steps through the support conclusion use only its
$L^p$ boundedness of $\A_k$ and the Abel theorem \eqref{eq:abel}.  The upper
bound on $p$ first enters in the next sentence, when the authors invoke a
theorem forcing an $L^p$ function with Fourier support on a sphere to vanish.

\section{Proof Intuition}

For one radius, a high-$p$ spectral component may survive because its central
Fourier transform can live on a Laguerre zero-sphere.  A second radius forces
the same component onto a second family of zero-spheres.  Unless the ratio of
the radii is a ratio of two zeros at the same Laguerre level, the two families
do not intersect.  The allowed Fourier support is then empty, so no
integrability-based support theorem is needed.  Abel summability reconstructs
$f$ from its now-vanishing spectral components.  At the exceptional ratios,
the source's explicit spherical eigenfunctions show that this obstruction is
real and give common kernel vectors.

\section{Exact finite-p classification}

\begin{theorem}\label{thm:main}
Let $G$, $m$, $n$, and $\mu_r$ be as above, and let $r,s>0$.
\begin{enumerate}
\item If $1\leq p\leq 2m/(m-1)$, then $f*\mu_r=0$ already implies
      $f=0$ for $f\in L^p(G)$.
\item If $2m/(m-1)<p<\infty$, then
\[
 f*\mu_r=f*\mu_s=0\quad\Longrightarrow\quad f=0
\]
for every $f\in L^p(G)$ if and only if $r/s\notin\E_n$.
\item If $p=\infty$ and $r/s\in\E_n$, the pair is not injective.  The
      nonexceptional $p=\infty$ case is not decided here.
\end{enumerate}
Consequently, the two-radius problem is completely classified on every finite
$L^p(G)$.
\end{theorem}

\begin{proof}
Part (1) is source Theorem~1.1 \cite{N_SY}.  We prove the new high-$p$
classification.

Assume first that $r/s\notin\E_n$, that
$2m/(m-1)<p<\infty$, and that $f*\mu_r=f*\mu_s=0$.  Since $m\geq2$, the
inequality $p>2m/(m-1)$ implies $p>2$.  Lemma~\ref{lem:support}, applied at
both radii, gives
\[
 \operatorname{supp}\mathcal F_t(\A_kf(z,\cdot))
 \subset Z_{k,r}\cap Z_{k,s}
\]
for every $k$ and almost every $z$.

For $k=0$, $L_0^{n-1}=1$, hence both zero sets are empty.  Let $k\geq1$.
If $a$ belonged to their intersection and $\lambda=|a|$, then
$\lambda>0$ because $L_k^{n-1}(0)\neq0$.  For some $i,j$ we would have
\[
 \frac{\lambda r^2}{2}=\zeta_{k,i},
 \qquad
 \frac{\lambda s^2}{2}=\zeta_{k,j}.
\]
Dividing yields
$r/s=\sqrt{\zeta_{k,i}/\zeta_{k,j}}\in\E_n$, a contradiction.  Thus every
support intersection is empty.  A distribution with empty support is zero,
so $\A_kf=0$ for all $k$.  Equation~\eqref{eq:abel} now gives $f=0$.

Conversely, suppose $r/s\in\E_n$.  Choose $k\geq1$ and roots
$\zeta_{k,i},\zeta_{k,j}$ such that
$r/s=\sqrt{\zeta_{k,i}/\zeta_{k,j}}$, and put
\[
 \lambda=\frac{2\zeta_{k,i}}{r^2}
        =\frac{2\zeta_{k,j}}{s^2}.
\]
The source constructs the nonzero spherical eigenfunction
\[
 F_{k,\lambda}(z,t)=
 \frac{J_{m/2-1}(\lambda|t|)}
      {(\lambda|t|)^{m/2-1}}\,\varphi_k^\lambda(z),
\]
up to a harmless normalization, and proves
\[
 F_{k,\lambda}*\mu_u
   =c_{k,n}\varphi_k^\lambda(u)F_{k,\lambda}.
\]
Our choice of $\lambda$ makes both
$\varphi_k^\lambda(r)$ and $\varphi_k^\lambda(s)$ zero.  Hence both means
annihilate $F_{k,\lambda}$.  The source also proves
$F_{k,\lambda}\in L^p(G)$ exactly for
$p>2m/(m-1)$, including $p=\infty$.  This proves the converse in (2) and the
noninjectivity assertion in (3).
\end{proof}

\section{Limitations and novelty status}

The theorem does not settle injectivity on $L^\infty(G)$ for
$r/s\notin\E_n$.  The finite-$p$ proof ends with norm-convergent Abel
reconstruction, which the source does not provide at the endpoint.  The
second source question about averages over general $K$-orbits is also not
addressed.

The local run indexes contained no result for arXiv:2109.14963.  Two bounded
primary-source searches used the exact title/question, authors, and close
variants involving ``two radius'', ``$H$-type'', spherical means, and Laguerre
zeros.  They found the source and its 2024 publication, and narrower adjacent
results, but no statement of the exact all-finite-$p$ classification above.
Novelty confidence is therefore moderate rather than definitive.

\section{Human review recommendation}

Send to a specialist in harmonic analysis on nilpotent Lie groups.  The main
audit point is Lemma~\ref{lem:support}: verify directly that the source's
support-containment argument remains valid for every finite $p>2$, and that
the source's upper bound on $p$ is used only after that containment has been
obtained.  A separate citation search under Wiener--Tauberian and spectral
synthesis terminology is also recommended.

\begin{thebibliography}{9}
\bibitem{N_SY}
E. K. Narayanan, P. K. Sanjay, and K. T. Yasser,
\emph{Injectivity of spherical means on $H$-type groups},
arXiv:2109.14963; J. Fourier Anal. Appl. \textbf{30} (2024), Article 33,
doi:10.1007/s00041-024-10089-9.
\end{thebibliography}

\end{document}
